\documentclass[oneside,12pt,a4paper]{book}

% [ Packages ]
\usepackage{geometry}
\usepackage{graphicx}
\usepackage{hyperref}

% UTF8 für Umlaute
\usepackage[utf8]{inputenc}

% Formale Zitate
\usepackage{csquotes}

% Neue deutsche Rechtschreibung
\usepackage[ngerman]{babel}

% Zitate unter Überschriften darstellen
\usepackage{epigraph}

% Zeilenabstand
\usepackage{setspace}

% Akronyme
\usepackage[
  printonlyused
]{acronym}

% Kapitel / Überschriften und Seitenzahlen im Kopf anzeigen
\usepackage{fancyhdr}

% BibLaTeX mit Biber
\usepackage[
  style=apa
  ,backend=biber
  ,apabackref=true
  ,language=auto
  ,abbreviate=true
  ,maxcitenames=2
  ,natbib=true
  ,hyperref=true
  ,maxbibnames=99
  ,url=true
]{biblatex}

% Sprache für APA zuweisen
\DeclareLanguageMapping{ngerman}{ngerman-apa}

% [ Bibliography ]
\bibliography{bib/literatur}

% [ Konfiguration ]
% [File: config.tex]
% Anpassung des Seitenlayouts
\geometry{
  a4paper
  ,top=29mm
  ,left=29mm
  ,right=29mm
  ,bottom=24mm
  ,headheight=15pt
% ,showframe
}

% Zitat nach Kapiteln -- Entferne Linie
\setlength{\epigraphrule}{0pt}

% [ Benutzerangaben ]
% [ Angaben zur Person ]
% Vorname
\newcommand{\DocumentAuthorPrename}{[Max]}
% Nachname
\newcommand{\DocumentAuthorName}{[Mustermann]}
% Straße und Hausnummer
\newcommand{\DocumentAuthorStreet}{[Musterstraße 67]}
% PLZ
\newcommand{\DocumentAuthorZip}{[123456]}
% Ort
\newcommand{\DocumentAuthorCity}{[Musterstadt]}
% E-Mail-Adresse
\newcommand{\DocumentAuthorEmail}{[vorname.nachname\(@\)hs-duesseldorf.de]}
% Studiengang
\newcommand{\DocumentThesisCourse}{[Medieninformatik/Medientechnik]}
% Matrikel-Nr.
\newcommand{\DocumentThesisMatr}{[123456]}

% [ Angaben zur Arbeit ]
% Art der Abschlussarbeit
\newcommand{\DocumentThesisType}{Wissenschaftliche Vertiefung und Bachelorarbeit}
% Titel der Abschlussarbeit
\newcommand{\DocumentThesisTitle}{[Titel der Wiss. Vertiefung und Bachelorarbeit]}
% Untertitel der Abschlussarbeit
\newcommand{\DocumentThesisSubtitle}{[Untertitel der Wiss. Vertiefung und Bachelorarbeit]}
% Zu erlangender akademischer Grad
\newcommand{\DocumentThesisDegree}{[Bachelor of Science]}
% Datum auf Cover
\newcommand{\DocumentThesisDateCover}{[Januar 1970]}
% Erstprüfer
\newcommand{\DocumentThesisFirstExaminer}{Prof. Dr.-Ing. Holger Schmidt}
% Zweitprüfer
\newcommand{\DocumentThesisSecondExaminer}{[Prof. Dr. Maxi Musterfrau]}

\begin{document}

  % Cover
  \thispagestyle{empty}

% Logo
\begin{figure}
  % Logo HSD
  \flushleft % linksbündig
  \includegraphics[width=.5\textwidth]{figures/hsd_m_logo.pdf}
\end{figure}

% Abstand
\vspace*{12ex}

\hfill
\begin{minipage}{.8\textwidth}
  \begin{flushright}
    
    {\Huge\bfseries\DocumentThesisTitle\par}
    
    \vspace*{2ex}		
    
    {\large\DocumentThesisSubtitle}
    
    \vspace*{12ex}
    
    \onehalfspacing
    
    \textbf{\DocumentThesisType}\\
    \vspace*{1ex}
    im Studiengang \DocumentThesisCourse\\
    zur Erlangung des akademischen Grades\\
    \DocumentThesisDegree\\
    
    \vspace*{8ex}
    
    Fachbereich Medien\\
    Hochschule Düsseldorf
    
    \vspace*{12ex}
    
    \DocumentAuthorPrename\ \DocumentAuthorName\\
    Matrikel-Nr.: \DocumentThesisMatr\\
    Datum: \DocumentThesisDateCover
    
    \vspace*{12ex}
    
    Erst- und Zweitprüfer\\
    \vspace*{1ex}
    \DocumentThesisFirstExaminer\\
    \DocumentThesisSecondExaminer
    
  \end{flushright}
\end{minipage}

  % Seitenzahlen röm.
  \pagenumbering{roman}

  % Eidesstattliche Erklärung
  % [ File: eidesstattliche_erklaerung.tex ]

\chapter*{Eidesstattliche Erklärung}

Ich erkläre hiermit an Eides statt, dass ich die vorliegende \DocumentThesisType\ selbständig und ohne unzulässige fremde Hilfe angefertigt habe. Die verwendeten Quellen sind vollständig zitiert. Diese Arbeit wurde weder in gleicher noch ähnlicher Form einem anderen Prüfungsamt vorgelegt oder veröffentlicht. Ich erkläre mich ausdrücklich damit einverstanden, dass diese Arbeit mittels eines Dienstes zur Erkennung von Plagiaten überprüft wird.

\vspace{2cm}

\begin{tabbing}
  \hspace{6cm} \= \kill
  \textit{Ort, Datum} \> \textit{\DocumentAuthorPrename\ \DocumentAuthorName}
\end{tabbing}

\vspace{8,5cm}

\noindent
\textbf{Kontaktinformationen}\\
\DocumentAuthorPrename\ \DocumentAuthorName\\
\DocumentAuthorStreet\\
\DocumentAuthorZip\ \DocumentAuthorCity\\
---\\
\DocumentAuthorEmail
\indent

  % Abstract
  \chapter*{Zusammenfassung}
\begin{center}
  \DocumentThesisTitle
  \\
  \vspace{.3cm}
  \DocumentAuthorPrename\ \DocumentAuthorName
\end{center}
Dies ist die Zusammenfassung Ihrer Arbeit \dots

\newpage

\chapter*{Abstract}
\begin{center}
  \DocumentThesisTitle
  \\
  \vspace{.3cm}
  \DocumentAuthorPrename\ \DocumentAuthorName
\end{center}
Diese möchten Sie natürlich auch auf Englisch bereitstellen \dots




  % Inhaltsverzeichnis

  \pagestyle{plain}
  \tableofcontents

  % Abbildungsverzeichnis
  \listoffigures

  % Tabellenverzeichnis
  \listoftables

  % Abkürzungsverzeichnis
  \newpage
  % [ File: acronyms.tex ]

% Liste der Abkürzungen
\newpage
\chapter*{Abkürzungsverzeichnis}
\begin{acronym}
  \setlength{\itemsep}{-0.1cm}
  \acro{BSI}{Bundesamt für Sicherheit in der Informationstechnik}
  \acro{CISO}{Chief Information Security Officer}
  \acro{SLA}{Service Level Agreement}
\end{acronym}

  \newpage

  % Seitenzahlen arab.
  \pagenumbering{arabic}
  \pagestyle{fancy}
  
  \part{Wissenschaftliche Vertiefung}
  
  \part{Bachelorarbeit}

  % Kapitel 1 -- Einleitung
  % Kapitel
\chapter{Einleitung}

\epigraph{\glqq The user's going to pick dancing pigs over security every time.\grqq\bigskip}
{\textsc{Bruce Schneier}\\ ($\ast$1963)}

\noindent
Das Verfassen einer eigenständigen Bachelor- bzw. Masterarbeit \dots
\indent

\section{Hinweise}

Bitte lassen Sie ein Exemplar des Anmeldeformulars, welches Sie bei der Anmeldung Ihrer Arbeit im Prüfungsamt haben unterschreiben lassen und auf dem das Abgabedatum vermerkt ist, als zweite Seite dieses Dokumentes einbinden.

Bei Abgabe zeigen Sie Ihre drei Exemplare im Studienbüro vor und lassen die Abgabe auf den entsprechenden Formularen eintragen. Zwei Exemplare müssen nun persönlich den beiden Prüfern übergeben werden. Das dritte Exemplar ist für Sie bestimmt.
Bitte denken Sie auch daran, dass Sie die eidesstattliche Erklärung vor Abgabe unterschreiben.

\section{Struktur}

Die Arbeit sollte u.\,a. nachfolgende Inhalte berücksichtigen:

\begin{itemize}
  \item Titelseite
  \item Eidesstattliche Erklärung
  \item Zusammenfassung und Abstract (Englisch)
  \item Inhaltsverzeichnis, Abbildungsverzeichnis, Tabellenverzeichnis, Abkürzungsverzeichnis und Literaturverzeichnis
\end{itemize}

\noindent
Der Zitierstil sollte nach APA (American Psychological Asscociation) Style (\url{http://www.apastyle.org/}) erfolgen.

\section{Organisatorisches}

\begin{itemize}
 \item Es gilt die jeweils aktuelle Pr\"ufungsordnung (\S 15 in BMI PO vom 04.08.2010 bzw. \S 15 in MMI PO vom 16.06.2011).
 \item Abzugeben gebunden als Ausdruck und elektronisch als PDF
\end{itemize}

\section{Bewertungskriterien}

Die Bewertung einer Arbeit erfolgt unter anderem auf Grundlage von \textbf{Schwierigkeitsgrad}, \textbf{wissenschaftlicher Arbeitstechnik}, \textbf{ingenieurmäßiger Vorgehensweise}, \textbf{Stil} und \textbf{Form}.

\noindent
Das zugehörige Kolloquium wird vor allem basierend auf der \textbf{Wiedergabe der Inhalte}, der \textbf{Foliengestaltung}, \textbf{Stil} und \textbf{Form} bewertet.

  % Kapitel 2 -- Stile
  % Kapitel
\chapter{Stile}

\epigraph{\glqq The wise know their weakness too well to assume infallibility; and he who knows most, knows best how little he knows.\grqq\bigskip}%\textsc
{{Thomas Jefferson}\\ (1743--1826)}

\noindent
Nachfolgend sind einige Beispiele zum Styling von Inhalten aufgeführt.
Eine gute Einführung in das Arbeiten mit \LaTeX\ bietet die Ausarbeitung von Jürgens und Feuerstack der FernUniversität in Hagen: \url{https://www.fernuni-hagen.de/imperia/md/content/zmi_2010/a026_latex_einf.pdf}.

\indent

\section{Text}

Dies ist ein Beispiel für \textit{kursiven} und \textbf{fetten} Text.

\bigskip

Abkürzungen werden in der Datei \texttt{acronyms.tex} definiert und können dann vereinfacht genutzt werden. Alle tatsächlich eingesetzten Abkürzungen werden automatisch im Abkürzungsverzeichnis aufgeführt. Eine Abkürzung wird bei der ersten Verwendung zusätzlich ausgeschrieben dargestellt. Ein Beispiel: Das \ac{BSI} stellte fest \dots

\section{Abbildungen und Tabellen}

Eine einfache Abbildung (\ref{fig:hsd_logo}):

\begin{figure}[hbtp]
  \centering
  \includegraphics[width=.6\textwidth]{figures/hsd_m_logo.pdf}
  \caption{Logo Hochschule Düsseldorf}
  \label{fig:hsd_logo}
\end{figure}

\bigskip

Eine einfache Tabelle (\ref{tab:beispieltabelle}):

\begin{table}[hbtp]
  \begin{center}
    \begin{tabular}{|c|c|}
      \hline 
      \rule[-1ex]{0pt}{2.5ex} Eins & 1 \\ 
      \hline 
      \rule[-1ex]{0pt}{2.5ex} Zwei & 2 \\ 
      \hline
    \end{tabular}
  \end{center}
  \caption{Beispieltabelle}
  \label{tab:beispieltabelle}
\end{table}

\section{Zitieren}

Die benötigte Literatur wird in der Datei \texttt{literatur.bib} gepflegt. Das Literaturverzeichnis wird automatisch generiert.

Dies ist ein Zitat von \parencite{iso27001:2013} \dots laut \parencite[65\psqq]{iso27001:2013} ist dieses Vorgehen empfehlenswert.

\section{Listen}

Unsortierte Liste:

\begin{itemize}
 \item Eins
 \item Zwei
 \item Drei
\end{itemize}

\bigskip

\noindent
Nummerierte Liste:

\begin{enumerate}
 \item Element
 \item Element
 \item Element
\end{enumerate}

\indent

  % Kapitel 3 -- Tooling
  % Kapitel
\chapter{Tools}

\epigraph{\glqq Man is still the most extraordinary computer of all.\grqq\bigskip}%\textsc
{{John F. Kennedy}\\ (1917--1963)}

\noindent
Nachfolgende Hinweise und Empfehlungen zum Einsatz von Tools vereinfachen den Umgang mit \LaTeX\ und Git.
\indent

\section{\LaTeX}

Das Verfassen von Dokumenten mit \LaTeX\ kann durch unterschiedlichste Tools unterstützt werden. Da \LaTeX\ grundsätzlich textbasiert arbeitet können jegliche Inhalte auch in einem einfachen Texteditor erstellt und angepasst werden.\\

Mittels unterschiedlichster Editoren kann die Erstellung und Pflege von Dokumenten mit \LaTeX\ vereinfacht werden. Unter Linux bietet der Editor „Kile“ (\url{http://kile.sourceforge.net/}) eine Vielzahl nützlicher Funktionen. Für Apple OS X und Microsoft Windows ist „Texmaker“ (\url{http://www.xm1math.net/texmaker/}) empfehlenswert.
  
    % Kapitel 4 -- Infrastruktur
  % Kapitel
\chapter{Infrastruktur}
\label{chap:infrastruktur}

\epigraph{\glqq The secret of all victory lies in the organization of the non-obvious.\grqq\bigskip}%\textsc
{{Marcus Aurelius}\\ (121--180)}

\noindent
Die Erstellung einer Arbeit sollte in einer bereitgestellten Infrastruktur erfolgen, die insbesondere bei der Planung und Verwaltung einer Arbeit unterstützt.
\indent

\section{GitLab-Server}

Unterschiedliche Dienste, die im Kontext einer Arbeit von Nutzen sind, werden über einen GitLab-Server bereitgestellt. Jeder Kandidat erhält einen persönlichen Zugang und ein eigenes Repository. In diesem Repository werden Dokumente und eigene Inhalte der Arbeit zentral verwaltet und somit dem Betreuer zur Kontrolle übergeben. Der GitLab-Server stellt hierzu in erster Linie ein Repository bereit. Das Repository bzw. eine Versionsverwaltung im Allgemeinen hilft vor allem bei der Verwaltung von textbasierten Dateien, so z.\,B. Quellcode oder Dokumente in \LaTeX .\\

Zur Planung einer Arbeit und Kontrolle des Fortschritts erfolgt das Projektmanagement digital innerhalb von GitLab. Hierzu werden Milestones und Issues angelegt und während des Projektes gepflegt bzw. Fortschritte kontrolliert. Eine möglichst präzise Projektplanung hilft bei der Vermeidung von etwaigen zeitlichen Engpässen im Laufe der Erstellung einer Arbeit.

\section{Versionsverwaltung mit Git}

Die Verwaltung der Arbeit, die mittels \LaTeX\ verfasst wird, und aller zugehörigen Dateien bzw. Dokumente kann auf einfache und sehr transparente Weise mittels einer Versionsverwaltung erfolgen. Als Versionsverwaltung wird Git eingesetzt. Git kann über \url{https://git-scm.com/} heruntergeladen und installiert werden. Git steht für alle gängigen Betriebssysteme bereit.\\

Jegliche Änderungen und Ergänzungen werden von Git erkannt und aufgezeichnet. Erfolgte Änderungen sollten mittels sog. „Commits“ eingepflegt und beschrieben werden. Die Versionsverwaltung erfolgt in erster Linie auf dem lokalen System. Erfolgte Änderungen bzw. Fortschritte sollten -- nicht nur als Backup -- regelmäßig über den bereitgestellten GitLab-Server dem Betreuer zur Verfügung gestellt werden.

Die Arbeit mit Git kann sowohl auf der Kommandozeile als auch in Applikation mit UI erfolgen. Die Applikation „SourceTree“ ermöglicht beispielsweise die komfortable Verwaltung von Git-Repositories. Siehe auch: \url{https://www.sourcetreeapp.com/}\\

Grundsätzliche Tipps zum Umgang mit Git liefern die offizielle Dokumentation (\url{https://git-scm.com/doc}) und das „Git Cheat Sheet“ (\url{https://www.git-tower.com/blog/git-cheat-sheet/}).

  % Anhang A
  % 

  % Literaturverzeichnis
  \printbibliography

\end{document}